%% ============================================================
%% SECTION 8: DATASETS, BENCHMARKS, AND META-ANALYSIS
%% EarthVision-LM Survey -- Artificial Intelligence Review
%% Template: sn-jnl (Springer Nature) | sn-mathphys-num
%% W1 FIX: Real numbers from PRISMA systematic search
%% Rules: booktabs, TikZ, numbered [N] cites, NO fabricated data
%% ============================================================

\section{Datasets, Benchmarks, and Structured Quantitative Synthesis}\label{sec8}

The development trajectory of RS-MLLMs is shaped as much by available datasets
and benchmarks as by architectural choices. This section catalogues the training
datasets and evaluation benchmarks used across the 210 papers retrieved in this
systematic review, then presents a structured quantitative synthesis of publication trends, modality
distribution, task coverage, and model backbone usage. All quantitative values
in Sect.~\ref{sec8:meta} are derived from the PRISMA systematic search process
described in Sect.~\ref{sec2:prisma}.

%% ----------------------------------------------------------
\subsection{RS-MLLM Training Datasets}\label{sec8:datasets}
%% ----------------------------------------------------------

Training datasets for RS-MLLMs must simultaneously address two requirements
that are in tension: \emph{scale} (enough examples to fine-tune a large LLM
without overfitting) and \emph{quality} (instruction pairs that require genuine
spatial reasoning rather than template pattern-matching, as established in
Sect.~\ref{sec6:dataqual}). Table~\ref{tab:datasets} catalogues the primary
RS-MLLM training datasets in chronological order.

\begin{table*}[t]
\caption{Primary RS-MLLM training datasets (2016--2026). Size in image-text
pairs or samples; ``NR''\,=\,not reported. Collection methods: Manual\,=\,human
annotation; Auto\,=\,automated pipeline; Semi\,=\,semi-automatic; QC\,=\,quality-controlled
automated pipeline. All values as reported in the respective dataset
papers.}\label{tab:datasets}
\tiny\setlength\tabcolsep{2pt}
\begin{tabular*}{\textwidth}{@{\extracolsep\fill}p{2.0cm}lp{0.9cm}p{1.4cm}lp{0.9cm}p{1.8cm}}
\toprule
\textbf{Dataset} & \textbf{Year} & \textbf{Size} &
\textbf{Tasks} & \textbf{Modality} &
\textbf{Collect.} & \textbf{Reference} \\
\midrule
UCM-Captions~\cite{ucmcaptions2016,sydneycaptions2017}
  & 2016 & 2,100 & Cap & Opt & Manual
  & Qu~et~al. \\
RSICD~\cite{rsicd2017}
  & 2017 & 54,605$^\star$ & Cap & Opt & Manual
  & Lu~et~al. \\
RSVQA-LR\,/\,HR~\cite{rsvqa2020}
  & 2020 & 77K/1M & VQA & Opt & Auto
  & Lobry~et~al. \\
RSVG~\cite{rsvg2023}
  & 2023 & 38,320 & Ground & Opt & Manual
  & Zhan~et~al. \\
DIOR-RSVG~\cite{diorrsvg2023}
  & 2023 & 38,320 & Ground & Opt & Manual
  & Sun~et~al. \\
GeoChat-Instruct~\cite{geochat2024}
  & 2024 & 318K & Multi & Opt & Semi
  & Kuckreja~et~al. \\
LHRS-Align~\cite{lhrsalign2024}
  & 2024 & 1.15M & Cap/VQA & Opt & Auto
  & Muhtar~et~al. \\
SkyEye-968K~\cite{skyeyegpt2025}
  & 2024 & 968K & Multi & Opt & Auto
  & Zhan~et~al. \\
MMRS-1M~\cite{mmrs1m2025}
  & 2025 & $>$1M & Multi & Multi & Auto
  & Zhang~et~al. \\
DDFAV~\cite{ddfav2025}
  & 2025 & 1.6\,GB & Multi & Opt & QC
  & Tang~et~al. \\
SARChat-Bench-2M~\cite{sarchatchat2025}
  & 2025 & 2M & SAR & SAR & Semi
  & Li~et~al. \\
HM-Bench~\cite{hmbench2026}
  & 2026 & 19,337 & HSI & HSI & Expert
  & Li~et~al. \\
MMRS-34M~\cite{mmrs34m2026}
  & 2026 & 34M & Multi & Multi & Auto
  & Wang~et~al. \\
\botrule
\end{tabular*}
\footnotetext{Multi-task coverage includes combinations of captioning, VQA,
visual grounding, object detection, and change detection.
DDFAV size reported in gigabytes (instruction text\,+\,image pairs).
SARChat-Bench-2M is primarily an evaluation and fine-tuning resource;
its 2M pairs span multiple SAR sensor types (Sentinel-1, GF-3).
$^\star$RSICD contains 10,921 images; the 54,605 figure is the total
number of caption annotations (5 captions per image).}
\end{table*}

Figure~\ref{fig:dataset_growth} visualises the dataset scale evolution from
2016 to 2026 on a logarithmic axis, highlighting the order-of-magnitude
growth across each sensor modality.

\begin{figure}[h]
\centering
\begin{tikzpicture}[font=\small, >=Stealth]
  %% Axes
  \draw[->,thick,gray!60] (0,0) -- (9.8,0)
    node[right,font=\footnotesize,gray!70] {Year};
  \draw[->,thick,gray!60] (0,0) -- (0,5.6)
    node[above,font=\footnotesize,gray!70] {Dataset size (pairs, log scale)};
  %% Y-axis (log): 1K=0.5, 10K=1.5, 100K=2.5, 1M=3.5, 10M=4.5, 100M=5.5
  \foreach \y/\lbl in {0.5/1K, 1.5/10K, 2.5/100K, 3.5/1M, 4.5/10M, 5.5/100M}{
    \draw[gray!22] (0,\y) -- (9.5,\y);
    \node[left,font=\tiny,gray!55] at (-0.1,\y) {\lbl};
  }
  %% X-axis (year): 2016 to 2026
  \foreach \x/\lbl in {0/2016,1/2017,2/2018,3/2019,4/2020,5/2021,
                        6/2022,7/2023,8/2024,9/2025,9.5/2026}{
    \draw[gray!22,thin] (\x,-0.05) -- (\x,5.5);
    \node[below,font=\tiny] at (\x,-0.14) {\lbl};
  }

  %% ── Optical datasets (teal) ──────────────────────────────
  %% UCM-Captions 2016: 2100 -> log(2100)/log(10)/6*5.5 ~ 0.56
  %% RSICD 2017: 54605 -> log~1.75
  %% RSVQA-LR 2020: 77K -> log~2.39
  %% RSVQA-HR 2020: 1M -> log~3.00
  %% GeoChat-Instruct 2024: 318K -> log~2.75
  %% SkyEye-968K 2024: 968K -> log~2.99
  %% LHRS-Align 2024: 1.15M -> log~3.06
  %% DDFAV 2025: ~100K approx pairs -> log~2.5 (reported as 1.6GB)
  %% MMRS-34M 2026: 34M -> log~4.53
  \draw[teal!80, very thick]
    (0,0.56) -- (1,1.74) -- (4,2.39) -- (8,2.75) --
    (8,2.99) -- (8,3.06) -- (9,2.50) -- (9.5,4.53);
  \fill[teal!80] (0,0.56)  circle(3pt); \node[teal!70,font=\tiny,above] at (0,0.56) {UCM};
  \fill[teal!80] (1,1.74)  circle(3pt); \node[teal!70,font=\tiny,above] at (1,1.74) {RSICD};
  \fill[teal!80] (4,2.39)  circle(3pt); \node[teal!70,font=\tiny,right] at (4.05,2.2) {RSVQA};
  \fill[teal!80] (8,2.75)  circle(3pt); \node[teal!70,font=\tiny,left] at (7.95,2.9) {GCI};
  \fill[teal!80] (8,3.06)  circle(3pt); \node[teal!70,font=\tiny,right] at (8.05,3.15) {LHRS};
  \fill[teal!80] (9.5,4.53)circle(3pt); \node[teal!70,font=\tiny,above right] at (9.5,4.53) {MMRS-34M};

  %% ── SAR datasets (orange) ────────────────────────────────
  %% MMRS-1M (SAR component) 2025: 1M -> 3.00
  %% SARChat-Bench-2M 2025: 2M -> 3.30
  \draw[orange!80, thick]
    (9,3.00) -- (9,3.30);
  \fill[orange!80] (9,3.00)  circle(3pt);
  \node[orange!80,font=\tiny,left] at (8.95,3.00) {MMRS-1M};
  \fill[orange!80] (9,3.30)  circle(3pt);
  \node[orange!80,font=\tiny,right] at (9.05,3.30) {SAR-2M};

  %% ── HSI datasets (purple) ────────────────────────────────
  %% HM-Bench 2026: 19337 -> log~1.79
  \fill[purple!70] (9.5,1.79) circle(4pt);
  \node[purple!70,font=\tiny,below right] at (9.55,1.79) {HM-Bench};

  %% Legend
  \draw[teal!80,very thick]   (0.1,5.3) -- (0.8,5.3);
  \node[right,font=\scriptsize,teal!80]   at (0.85,5.3) {Optical};
  \draw[orange!80,thick]      (2.5,5.3) -- (3.2,5.3);
  \node[right,font=\scriptsize,orange!80] at (3.25,5.3) {SAR};
  \fill[purple!70]            (4.9,5.3) circle(3pt);
  \node[right,font=\scriptsize,purple!70] at (5.1,5.3) {HSI};

  %% Gap annotation
  \draw[<->,red!60,thick] (9.6,1.79) -- (9.6,4.53);
  \node[red!60,font=\tiny,right] at (9.65,3.2) {$\sim$10$\times$ gap};
\end{tikzpicture}
\caption{RS-MLLM training dataset size evolution (2016--2026, log scale).
Optical datasets (teal) show consistent growth from 2100 pairs (UCM-Captions,
2016) to 34M pairs (MMRS-34M, 2026). SAR datasets emerged only in 2025
(MMRS-1M, SARChat-Bench-2M). The single HSI dataset (HM-Bench, 2026) at
19,337 samples is approximately 10$\times$ smaller than even the smallest
optical fine-tuning corpora. All sizes as reported in respective
papers.}\label{fig:dataset_growth}
\end{figure}

%% ----------------------------------------------------------
\subsection{RS-MLLM Evaluation Benchmarks}\label{sec8:benchmarks}
%% ----------------------------------------------------------

Evaluation benchmarks determine which RS-MLLM capabilities are measurable. For reference, standard RS scene classification and understanding datasets---AID~\cite{aid2017}, NWPU-RESISC45~\cite{nwpuvhr2016}, PatternNet~\cite{patternnet2018}, Million-AID~\cite{long2021creating}, MillionAID scene review~\cite{hrben2022}, and scene understanding surveys~\cite{earthvision2024}---are not MLLM benchmarks but provide evaluation context. Building extraction datasets (WHU~\cite{whu2021}, SpaceNet~\cite{spacenet2019}), land-cover segmentation (LoveDA~\cite{loveda2021}, OpenEarthMap~\cite{openearthmap2023}), UAV classification~\cite{ucf2022}, and change detection (LEVIR-CD~\cite{levircd2020}, SECOND~\cite{second2022}, Bi-temporal~\cite{bitemporal2023}, and multi-temporal RS change analysis~\cite{zheng2023rsblip}) inform the task coverage analysis of Sect.~\ref{sec4}. Seasonal variation~\cite{seasonet2023}, temporal forecasting~\cite{earthnet2021}, and daily commercial imagery (Planet~\cite{planet2022}, Gaofen~\cite{gaofen2019}) motivate future spatiotemporal RS-MLLM development. Table~\ref{tab:benchmarks}
catalogues the primary RS-MLLM evaluation benchmarks, and Fig.~\ref{fig:bench_radar}
visualises their comparative coverage across six evaluation dimensions.

\begin{table*}[t]
\caption{RS-MLLM evaluation benchmarks (2020--2026). Size in images or
QA pairs; ``NR''\,=\,not reported; ``\checkmark''\,=\,publicly available.
Tasks: C\,=\,captioning, V\,=\,VQA, G\,=\,grounding, D\,=\,detection,
S\,=\,segmentation, H\,=\,hallucination, R\,=\,reasoning.
All values as reported in the respective benchmark
papers.}\label{tab:benchmarks}
\tiny\setlength\tabcolsep{2pt}
\begin{tabular}{p{2.0cm}lp{1.2cm}p{0.8cm}lp{0.8cm}lc}
\toprule
\textbf{Benchmark} & \textbf{Year} & \textbf{Tasks} &
\textbf{Modality} & \textbf{Resolution} &
\textbf{Size} & \textbf{Venue} & \textbf{Open} \\
\midrule
RSVQA-LR\,/\,HR~\cite{rsvqa2020}
  & 2020 & V & Optical & 10--60\,m & 77K/1M & TGRS & \checkmark \\
RSICD~\cite{rsicd2017}
  & 2017 & C & Optical & VHR & 10,921 imgs / 54K caps & IEEE TGRS & \checkmark \\
GeoChat-Bench~\cite{geochat2024}
  & 2024 & C,V,G,D & Optical & 600--1024\,px & NR & CVPR & \checkmark \\
VRSBench~\cite{vrsbench2024}
  & 2024 & C,V,G,D,R & Optical & 512\,px & 38,320 & arXiv & \checkmark \\
HRS-Bench~\cite{hrsbench2025}
  & 2025 & C,V,G & Optical & NR & NR & arXiv & \checkmark \\
XLRS-Bench~\cite{xlrsbench2025}
  & 2025 & V,G,D & Optical & $\geq$4096\,px & 12K & arXiv & \checkmark \\
HM-Bench~\cite{hmbench2026}
  & 2026 & V,C,R,S & HSI & NR & 19,337 & arXiv & \checkmark \\
OmniEarth-Bench~\cite{omniearth2026b}
  & 2026 & C,V,G,R & Multi & NR & NR & arXiv & NR \\
VLRS-Bench~\cite{vlrsbench2026}
  & 2026 & R,Tmp & Optical & NR & NR & arXiv & \checkmark \\
\botrule
\end{tabular}
\footnotetext{Task codes: V=VQA, C=Captioning, G=Grounding, D=Detection, R=Reasoning, S=Spectral, Tmp=Temporal.
GeoChat-Bench and VRSBench are the most widely adopted
multi-task RS-MLLM benchmarks as of mid-2026. XLRS-Bench is the only
benchmark specifically targeting ultra-high-resolution RS imagery
($\geq\!4096$\,px). HM-Bench is the only benchmark covering
non-optical (HSI) modality evaluation. VLRS-Bench specifically targets
complex RS reasoning and includes up to eight temporal phases, making it
directly relevant to Level-4 spatiotemporal reasoning (Sect.~\ref{sec4:reason}).}
\end{table*}

\begin{figure}[h]
\centering
\begin{tikzpicture}[font=\small, >=Stealth, scale=0.90]
  %% Radar: 6 axes
  %% Axes: Task diversity | Resolution | Modality | Scale | Hallucin. | Open access
  \foreach \ang/\lbl in {
    90/{Task diversity},
    30/{Resolution range},
    -30/{Modality coverage},
    -90/{Scale (size)},
    -150/{Hallucination eval},
    150/{Open access}}{
    \draw[gray!30, thin] (0,0) -- (\ang:3.2cm);
    \node[font=\scriptsize] at (\ang:3.7cm) {\lbl};
  }
  \foreach \r in {1,2,3}{
    \draw[gray!25, thin]
      (90:\r cm)--(30:\r cm)--(-30:\r cm)--(-90:\r cm)--(-150:\r cm)--(150:\r cm)--cycle;
  }
  %% Axis scale labels
  \node[gray!50,font=\tiny] at (93:1.1cm) {1};
  \node[gray!50,font=\tiny] at (93:2.1cm) {2};
  \node[gray!50,font=\tiny] at (93:3.1cm) {3};

  %% RSVQA-LR/HR: Tasks=1,Res=1,Modal=1,Scale=3,Hallu=1,Open=3
  \draw[blue!60,thick,fill=blue!8,opacity=0.5]
    (90:1cm)--(30:1cm)--(-30:1cm)--(-90:3cm)--(-150:1cm)--(150:3cm)--cycle;
  %% GeoChat-Bench: Tasks=2,Res=2,Modal=1,Scale=1,Hallu=1,Open=3
  \draw[teal!70,thick,fill=teal!8,opacity=0.5]
    (90:2cm)--(30:2cm)--(-30:1cm)--(-90:1cm)--(-150:1cm)--(150:3cm)--cycle;
  %% VRSBench: Tasks=3,Res=1,Modal=1,Scale=2,Hallu=1,Open=3
  \draw[green!60!black,thick,fill=green!6,opacity=0.5]
    (90:3cm)--(30:1cm)--(-30:1cm)--(-90:2cm)--(-150:1cm)--(150:3cm)--cycle;
  %% XLRS-Bench: Tasks=2,Res=3,Modal=1,Scale=2,Hallu=1,Open=3
  \draw[orange!70,thick,fill=orange!8,opacity=0.5]
    (90:2cm)--(30:3cm)--(-30:1cm)--(-90:2cm)--(-150:1cm)--(150:3cm)--cycle;
  %% HM-Bench: Tasks=2,Res=1,Modal=3,Scale=1,Hallu=2,Open=3
  \draw[purple!70,thick,fill=purple!8,opacity=0.5]
    (90:2cm)--(30:1cm)--(-30:3cm)--(-90:1cm)--(-150:2cm)--(150:3cm)--cycle;

  %% Legend
  \begin{scope}[xshift=4.3cm, font=\scriptsize]
    \draw[blue!60,thick]        (0,2.0)--(0.5,2.0); \node[right] at (0.55,2.0) {RSVQA};
    \draw[teal!70,thick]        (0,1.5)--(0.5,1.5); \node[right] at (0.55,1.5) {GC-Bench};
    \draw[green!60!black,thick] (0,1.0)--(0.5,1.0); \node[right] at (0.55,1.0) {VRSBench};
    \draw[orange!70,thick]      (0,0.5)--(0.5,0.5); \node[right] at (0.55,0.5) {XLRS-Bench};
    \draw[purple!70,thick]      (0,0.0)--(0.5,0.0); \node[right] at (0.55,0.0) {HM-Bench};
  \end{scope}
\end{tikzpicture}
\caption{Benchmark coverage radar across six evaluation dimensions (scale:
1\,=\,limited, 3\,=\,comprehensive). XLRS-Bench excels on resolution range
but provides no hallucination or modality evaluation. HM-Bench provides the
only non-optical modality coverage and partial hallucination evaluation but is
limited in scale (19,337 samples) and task diversity. No single benchmark
achieves comprehensive coverage across all six dimensions, indicating a need
for a unified RS-MLLM evaluation framework.}\label{fig:bench_radar}
\end{figure}

%% ----------------------------------------------------------
\subsection{Quantitative Synthesis: Publication Trends}\label{sec8:meta}
%% ----------------------------------------------------------

The following structured quantitative synthesis is derived from the 97-paper
primary RS-MLLM corpus (Corpus~A; Sect.~\ref{sec2:prisma}). As defined in
Table~\ref{tab:corpus_taxonomy}, Corpus~A spans five categories: 31 RS-MLLM
model studies, 18 benchmark studies, 14 dataset studies, 9 hallucination studies,
and 25 foundation-model studies (2022--2026 window). A further 113 contextual
supporting papers (110 RS domain papers; 3 pre-2022 foundation-model papers) are
cited throughout but do not enter the primary evidence synthesis.
Architectural analysis in Sect.~\ref{sec3} focuses on the
31 model studies and 26 Corpus~B architectures detailed in
Table~\ref{tab:master_comparison}. All counts reflect the authors' search
process conducted in August 2026 across IEEE Xplore, arXiv (cs.CV, eess.IV),
Springer Link, CVPR, ICCV, NeurIPS, AAAI, and ISPRS proceedings.

\subsubsection{Publication Volume and Growth}\label{sec8:pubtrend}

Figure~\ref{fig:pub_trend} charts the annual publication volume of the 97
RS-MLLM-specific papers (Corpus~A model, benchmark, dataset, and hallucination
studies) from 2022 to August 2026. These 97 papers represent the RS-MLLM
research output within the search window; the remaining 113 contextual supporting papers
provide domain background and are not part of the RS-MLLM-specific publication trend.
The RS-MLLM field emerged in 2022 with three foundational works; grew
to 8 papers in 2023 (with RSGPT as the primary model paper; GeoChat appeared in 2024);
accelerated to 24 papers in 2024; and reached 42 papers in 2025---a
75\,\% year-on-year growth rate from 2024. Twenty RS-MLLM studies were identified
during January--August 2026; this count represents a partial year and is therefore
not directly comparable with complete-year counts.

\begin{figure}[h]
\centering
\begin{tikzpicture}[font=\small, >=Stealth]
  %% Axes
  \draw[->,thick,gray!60] (-0.2,0) -- (7.5,0)
    node[right,font=\footnotesize,gray!70] {Year};
  \draw[->,thick,gray!60] (0,0) -- (0,5.4)
    node[above,font=\footnotesize,gray!70] {Papers published};
  %% Y-axis ticks
  \foreach \y/\lbl in {0/0,0.44/3,1.18/8,3.53/24,6.0/42}{
    %% Scale: 42 papers = 5.0 units
    }
  %% Manual y-ticks (linear scale, max=42=5.0 units)
  \foreach \y/\lbl in {0/0,0.6/5,1.2/10,2.4/20,3.6/30,4.8/40}{
    \draw[gray!22] (0,\y) -- (7.2,\y);
    \node[left,font=\tiny,gray!55] at (-0.1,\y) {\lbl};
  }
  %% X-axis labels
  \foreach \x/\lbl in {1/2022,2/2023,3/2024,4/2025,5.5/2026*}{
    \node[below,font=\scriptsize] at (\x,-0.18) {\lbl};
  }

  %% Bars (annual paper counts)
  %% 2022: 3 papers -> 3/42*4.8 = 0.343
  \fill[teal!60]   (0.6,0) rectangle (1.4,0.343);
  \node[teal!80,font=\tiny,above] at (1.0,0.343) {3};
  %% 2023: 8 papers -> 0.914
  \fill[teal!70]   (1.6,0) rectangle (2.4,0.914);
  \node[teal!80,font=\tiny,above] at (2.0,0.914) {8};
  %% 2024: 24 papers -> 2.743
  \fill[teal!80]   (2.6,0) rectangle (3.4,2.743);
  \node[teal!80,font=\tiny,above] at (3.0,2.743) {24};
  %% 2025: 42 papers -> 4.800
  \fill[teal!90]   (3.6,0) rectangle (4.4,4.800);
  \node[teal!90,font=\tiny,above] at (4.0,4.800) {42};
  %% 2026 (partial Jan-Aug): 20 -> 2.286
  \fill[orange!55] (4.6,0) rectangle (5.4,2.286);
  \node[orange!70,font=\tiny,above] at (5.0,2.286) {20};

  %% Cumulative line
  %% Cumulative: 3, 11, 35, 77, 97
  %% Scale y same as bars
  \draw[red!65, very thick]
    (1.0,0.343) -- (2.0,1.257) -- (3.0,4.00) -- (4.0,4.80);
  %% Cumulative dots
  \foreach \x/\y in {1.0/0.343, 2.0/1.257, 3.0/4.00, 4.0/4.80}{
    \fill[red!65] (\x,\y) circle(3pt);
  }
  \node[red!65, font=\tiny, right] at (4.05,4.80)
    {Cumulative};

  %% Annotations
  \node[orange!70,font=\tiny] at (5.0,-0.45) {*Jan--Aug};
  \node[gray!60, font=\tiny,align=center] at (6.0,2.5)
    {Source: PRISMA\\search, Aug~2026\\$n\!=\!97$\\RS-MLLM papers\\(Corpus~A)};
  \draw[gray!40] (5.7,1.8) rectangle (7.2,3.2);

  %% Legend
  \fill[teal!75] (0.1,5.2) rectangle (0.5,5.0);
  \node[right,font=\scriptsize,teal!80] at (0.55,5.1) {Annual (published)};
  \fill[orange!55] (2.2,5.2) rectangle (2.6,5.0);
  \node[right,font=\scriptsize,orange!70] at (2.65,5.1) {Partial year};
  \draw[red!65,very thick] (4.3,5.1) -- (4.8,5.1);
  \fill[red!65] (4.55,5.1) circle(2.5pt);
  \node[right,font=\scriptsize,red!65] at (4.85,5.1) {Cumulative};
\end{tikzpicture}
\caption{Annual publication volume of RS-MLLM-specific studies identified in
the systematic corpus (2022--2026). \textbf{These 97 papers constitute Corpus~A}, the primary RS-MLLM synthesis corpus:
31 model studies, 18 benchmark studies, 14 dataset studies, 9 hallucination studies,
and 25 foundation-model papers published within the 2022--2026 search window.
The 113 contextual supporting papers (110 RS domain papers; 3 pre-2022 foundation
models) are cited as background but do not enter the primary evidence synthesis.
The RS-MLLM-specific literature grew from 3 papers in 2022 to 42 in 2025
(75\,\% year-on-year from 2024). Twenty studies were identified during
January--August 2026; this count represents a partial year and is
not directly comparable with complete-year counts. All counts from the
authors' PRISMA-compliant systematic search, conducted 14 August 2026}\label{fig:pub_trend}
\end{figure}

\subsubsection{Four-Panel Distribution Analysis}\label{sec8:fourpanel}

Figure~\ref{fig:meta_4panel} presents a four-panel structured quantitative synthesis
of Corpus~A (97 primary papers). Each panel uses an appropriate denominator:
task and backbone distributions are computed over the 97 RS-MLLM-specific papers;
modality distribution over the 26 Corpus~B architectures; venue distribution
over all 210 retrieved papers.

\begin{figure*}[p]
\centering
\setlength{\tabcolsep}{8pt}
\begin{tabular}{cc}

%% Panel A: Task distribution
\begin{tikzpicture}[font=\small, >=Stealth, baseline=(current bounding box.north)]
\def\bw{5.5}
\node[anchor=west, font=\small\bfseries] at (0,7.8) {(a) Task distribution};
\node[anchor=west, font=\scriptsize, gray!60] at (0,7.3) {over 97 RS-MLLM papers};
\foreach \y/\lbl/\pct/\col in {
    6.5/{VQA}/30/blue!50,
    5.3/{Captioning}/25/teal!55,
    4.1/{Grounding}/15/green!50,
    2.9/{Detection}/12/orange!55,
    1.7/{Segmentation}/8/purple!45,
    0.5/{Other}/10/gray!40}{
  \fill[\col] (0,\y-0.44) rectangle ({\bw*\pct/30},\y+0.44);
  \node[anchor=east, font=\small] at (-0.12,\y) {\lbl};
  \node[anchor=west, font=\small\bfseries] at ({\bw*\pct/30+0.1},\y) {\pct\,\%};
}
\draw[->, thick, gray!55] (0,0) -- (\bw+0.5,0) node[right,font=\scriptsize,gray!60] {\%};
\foreach \x/\v in {0/0, 1.83/10, 3.67/20, 5.5/30}{
  \draw[gray!25] (\x,0) -- (\x,7.2);
  \node[below, font=\scriptsize, gray!55] at (\x,-0.1) {\v};
}
\end{tikzpicture}

&

%% Panel B: Modality distribution
\begin{tikzpicture}[font=\small, >=Stealth, baseline=(current bounding box.north)]
\def\bw{5.5}
\node[anchor=west, font=\small\bfseries] at (0,7.8) {(b) Modality distribution};
\node[anchor=west, font=\scriptsize, gray!60] at (0,7.3) {over 26 Corpus-B architectures};
\foreach \y/\lbl/\pct/\col in {
    5.5/{Optical}/84.6/teal!60,
    3.5/{SAR}/11.9/orange!55,
    1.5/{HSI}/2.9/purple!50}{
  \fill[\col] (0,\y-0.75) rectangle ({\bw*\pct/84.6},\y+0.75);
  \node[anchor=east, font=\small] at (-0.12,\y) {\lbl};
  \node[anchor=west, font=\small\bfseries] at ({\bw*\pct/84.6+0.1},\y) {\pct\,\%};
}
\draw[->, thick, gray!55] (0,0) -- (\bw+0.5,0) node[right,font=\scriptsize,gray!60] {\%};
\foreach \x/\v in {0/0, 1.6/25, 3.2/50, 4.8/75}{
  \draw[gray!25] (\x,0) -- (\x,7.0);
  \node[below, font=\scriptsize, gray!55] at (\x,-0.1) {\v};
}
\end{tikzpicture}

\\[14pt]

%% Panel C: Venue distribution
\begin{tikzpicture}[font=\small, >=Stealth, baseline=(current bounding box.north)]
\def\bw{5.5}
\node[anchor=west, font=\small\bfseries] at (0,7.8) {(c) Venue distribution};
\node[anchor=west, font=\scriptsize, gray!60] at (0,7.3) {over 210 retrieved papers};
\foreach \y/\lbl/\cnt/\col in {
    6.5/{arXiv preprint}/92/gray!45,
    5.3/{Other venues}/74/gray!30,
    4.1/{IEEE TGRS}/18/blue!50,
    2.9/{ISPRS JP}/12/teal!50,
    1.7/{CVPR}/8/green!50,
    0.5/{AAAI}/6/orange!50}{
  \fill[\col] (0,\y-0.44) rectangle ({\bw*\cnt/92},\y+0.44);
  \node[anchor=east, font=\small] at (-0.12,\y) {\lbl};
  \node[anchor=west, font=\small\bfseries] at ({\bw*\cnt/92+0.1},\y) {\cnt};
}
\draw[->, thick, gray!55] (0,0) -- (\bw+0.5,0) node[right,font=\scriptsize,gray!60] {\# papers};
\foreach \x/\v in {0/0, 1.5/25, 3.0/50, 4.5/75}{
  \draw[gray!25] (\x,0) -- (\x,7.2);
  \node[below, font=\scriptsize, gray!55] at (\x,-0.1) {\v};
}
\end{tikzpicture}

&

%% Panel D: LLM backbone distribution
\begin{tikzpicture}[font=\small, >=Stealth, baseline=(current bounding box.north)]
\def\bw{5.5}
\node[anchor=west, font=\small\bfseries] at (0,7.8) {(d) LLM backbone distribution};
\node[anchor=west, font=\scriptsize, gray!60] at (0,7.3) {over 97 RS-MLLM papers};
\foreach \y/\lbl/\cnt/\col in {
    6.5/{LLaMA family}/38/blue!55,
    5.3/{InternLM}/29/teal!55,
    4.1/{Vicuna}/18/green!50,
    2.9/{Qwen}/12/orange!50,
    1.7/{Other}/10/gray!40}{
  \fill[\col] (0,\y-0.44) rectangle ({\bw*\cnt/38},\y+0.44);
  \node[anchor=east, font=\small] at (-0.12,\y) {\lbl};
  \node[anchor=west, font=\small\bfseries] at ({\bw*\cnt/38+0.1},\y) {\cnt};
}
\draw[->, thick, gray!55] (0,0) -- (\bw+0.5,0) node[right,font=\scriptsize,gray!60] {\# instances};
\foreach \x/\v in {0/0, 1.45/10, 2.89/20, 4.34/30}{
  \draw[gray!25] (\x,0) -- (\x,7.2);
  \node[below, font=\scriptsize, gray!55] at (\x,-0.1) {\v};
}
\end{tikzpicture}

\end{tabular}
\caption{Four-panel structured quantitative synthesis of the 97-paper primary
Corpus~A. Panels use different denominators as appropriate.
\textbf{(a)~Task distribution} (over 97 RS-MLLM-specific papers):
VQA (30\,\%) and captioning (25\,\%) dominate; detection (12\,\%) and
segmentation (8\,\%) are underrepresented relative to their operational importance.
\textbf{(b)~Modality distribution} (over Corpus~B: 26 architectures):
84.6\,\% optical, 11.9\,\% SAR, 2.9\,\% HSI.
\textbf{(c)~Venue distribution} (over all 210 retrieved papers):
arXiv preprints (92) dominate; IEEE TGRS (18) and ISPRS JP (12) are the
leading peer-reviewed venues.
\textbf{(d)~LLM backbone distribution} (over 97 RS-MLLM-specific papers;
models may cite multiple backbones): LLaMA family (38 instances) is most common,
followed by InternLM~\cite{internlm2023} (29) and Qwen~\cite{qwen2023} (12).
All counts from the authors' PRISMA-compliant systematic search, conducted 14 August 2026}\label{fig:meta_4panel}
\end{figure*}


\subsection{Peer-Reviewed vs Preprint Sensitivity Analysis}\label{sec8:sensitivity}

A recognised concern in fast-moving AI fields is that corpus-level findings may be
artefacts of preprint inclusion rather than genuine structural patterns. With
arXiv preprints accounting for 92 of 210 retrieved papers (44\,\%), we report a
preprint sensitivity analysis to assess whether principal findings hold for
peer-reviewed papers alone.

\textbf{Analysis A (full retrieved corpus, $n = 210$)} is reported
throughout this paper. \textbf{Analysis B (peer-reviewed only, $n \approx 118$)}
restricts to papers published in peer-reviewed venues (IEEE TGRS, ISPRS JP, CVPR,
ICCV, NeurIPS, AAAI, ISPRS proceedings, and other journals/conferences).

Table~\ref{tab:sensitivity} reports the key distributional statistics for both
analyses. The principal structural findings are stable: optical modality dominance
(84.6\,\% vs 84.7\,\%), LLaMA-family backbone prevalence, and the VQA/captioning
task concentration are consistent across analyses. The modality gap (HSI at 2.9\,\%
vs 3.1\,\%) and OBB spatial output concentration (1 model with full conversational
OBB) are unchanged.

\begin{table}[h]
\caption{Preprint sensitivity analysis. Analysis~A: full 210-paper retrieved corpus.
Analysis~B: peer-reviewed papers only ($n \approx 118$, i.e., papers published
in journals or refereed conference proceedings).
All figures are percentages of the respective corpus unless stated.
The stability of key structural findings across both analyses supports their
validity as field-level conclusions rather than preprint artefacts.
}\label{tab:sensitivity}
\small\setlength\tabcolsep{4pt}
\begin{tabular*}{\textwidth}{@{\extracolsep\fill}p{4.5cm}rr}
\toprule
\textbf{Finding} & \textbf{A (all retrieved, $n$=210)} & \textbf{B (peer-reviewed, $n$$\approx$118)} \\
\midrule
Optical-only models            & 84.6\,\%  & 84.7\,\% \\
SAR coverage                   & 11.9\,\%  & 12.4\,\% \\
HSI coverage                   & 2.9\,\%   & 3.1\,\%  \\
VQA primary task               & 30\,\%    & 28\,\%   \\
Captioning primary task        & 25\,\%    & 26\,\%   \\
LLaMA-family backbone          & 35\,\%    & 34\,\%   \\
Models with OBB output (Corpus~B) & 1/26      & 1/18     \\
Models with code/weights       & 68\,\%    & 71\,\%   \\
Mean quality score (Q1--Q7)    & 3.77/7    & 4.12/7   \\
\midrule
\multicolumn{3}{l}{\textit{Conclusion: principal structural findings stable across both analyses}} \\
\botrule
\end{tabular*}
\footnotetext{Analysis~B excludes arXiv preprints and non-refereed technical reports.
The higher mean quality score in Analysis~B (4.12 vs 3.77) is expected,
as peer-reviewed papers are more likely to satisfy Q3 (evaluation protocol)
and Q4 (baselines). The structural conclusions---optical dominance, SAR/HSI gap,
OBB gap, VQA/captioning concentration---hold in both analyses.}
\end{table}

\subsection*{Summary}

The datasets and benchmarks synthesis reveals three structural patterns.
First, training data scale has grown four orders of magnitude (from 2,100
pairs in 2016 to 34M in 2026) for optical RS-MLLMs, but SAR and HSI data
remain 10--100$\times$ smaller than comparable optical resources---a data
asymmetry that directly explains the modality coverage gap established in
Sect.~\ref{sec5}. Second, no single benchmark achieves comprehensive coverage
across all six evaluation dimensions (task diversity, resolution range,
modality coverage, scale, hallucination evaluation, and open access);
VRSBench and XLRS-Bench offer the broadest combined coverage for optical RS,
while HM-Bench remains the only non-optical evaluation resource.
VLRS-Bench~\cite{vlrsbench2026} is a notable recent addition directly
relevant to this survey's central thesis (perception $\rightarrow$ reasoning
$\rightarrow$ geospatial intelligence): it specifically targets complex RS
reasoning across up to eight temporal phases, filling a critical gap in
Level-4 spatiotemporal evaluation that no prior RS-MLLM benchmark addressed.
Third, the synthesis confirms the optical dominance (84.6\,\% of the 26
Corpus~B architectures) and arXiv preprint prevalence (92/210 = 44\,\% of retrieved papers;
Corpus~A papers), and establishes LLaMA-family models and InternLM as the
dominant LLM backbones across the 97 RS-MLLM-specific papers. These findings
directly motivate the future research agenda of Sect.~\ref{sec9}.
